
The methodology is organized around four interconnected experimental questions, each with a distinct measurement protocol. The central design principle is oracle orbit construction: rather than inferring symmetry structure from natural variation in the zoo, we directly construct ground-truth functional equivalents by applying explicit symmetry transforms.

\subsubsection{3.1 Setting and Data}

\textbf{Model Zoo.} We use the Schürholt MNIST model zoo [Schürholt et al., 2022], a collection of 2-layer MLPs trained on MNIST with architecture $784\rightarrow 64\rightarrow 10$ (ReLU activations, no bias in the final layer). Models vary in learning rate (log-uniform in [$10^{-5}, 10^{-1}])$, weight decay, and random initialization seed. We use a locally archived subset of N=500 models with ground-truth property labels: test accuracy, generalization gap (train minus test accuracy), and learning rate recovery. Weights are stored as flat tensors of dimension D=51,850 ($784\times 64 + 64\times 10)$. The full Schürholt zoo contains approximately 50,000 models but was unavailable via HuggingFace at experiment time; all experiments use this 500-model local archive.

\subsubsection{3.2 Symmetry Orbit Construction}

\textbf{Scaling orbits.} For a 2-layer ReLU MLP with weight matrices $W_1 \in \mathbb{R}^{d_{\mathrm{in}} \times h}$ and $W_2 \in \mathbb{R}^{h \times d_{\mathrm{out}}}$, the scaling symmetry transform applies per-neuron positive rescaling: $W_{1}[$:,i] $\leftarrow \alpha _i W_{1}[$:,i] and $W_{2}[i$,:] $\leftarrow W_{2}[i$,:] / $\alpha _i$ for any $\alpha _i$ > 0. By positive homogeneity of ReLU (ReLU($\alpha z) = \alpha$ ReLU(z) for $\alpha$ > 0), the resulting network computes the same function as the original for all inputs. We sample scaling factors log-uniformly: log $\alpha _i$ ∼ Uniform(log 0.1, log 10), covering a $100\times$ dynamic range per neuron.

\textbf{Sign-flip orbits.} The sign-flip symmetry transform applies per-neuron sign changes: $W_{1}[$:,i] $\leftarrow$ s\_i $W_{1}[$:,i] and $W_{2}[i$,:] $\leftarrow$ s\_i $W_{2}[i$,:] for s\_i $\in$ {-1, +1}. For scaling ($\alpha _i$ > 0), the positive-homogeneity argument gives exact functional equivalence. For sign flips (s\_i = -1), however, ReLU is not odd: ReLU(-z) $\neq$ -ReLU(z) in general. The simultaneous two-layer flip is therefore an \textit{approximate} functional equivalence: the two networks agree on inputs where the neuron's pre-activation sign is unchanged by the flip, and disagree on inputs near the neuron's decision boundary. For gradient-trained networks with approximately symmetric weight distributions, this approximation holds well in practice. We treat sign-flip orbits as approximate functional equivalence classes and note that runtime functional equivalence verification on held-out inputs was not performed (see Section 6.2, Limitation 3). We sample signs uniformly: s\_i ∼ Bernoulli(0.5) mapped to {-1, +1}. We construct N=500 pairs per condition (A--E), yielding 2,500 oracle orbit pairs total across all five conditions.

\subsubsection{3.3 NFT Invariance Probe}

\textbf{Encoder.} We use NFT \citep{Zhou2023NeuralFunctionalNetworks} with d\_model=256, 4 self-attention layers, 8 heads, and CLS-token pooling. CLS pooling is critical --- mean pooling collapses embeddings and destroys model discrimination (verified empirically). NFT has approximately 3.4M parameters. NFT is trained on Condition A (raw weights) for property prediction and then frozen; embeddings are extracted without gradient computation.

\textbf{Invariance gap.} For each oracle orbit pair (b, b'), we extract embeddings e(b) and e(b'). Within-orbit similarity is cos(e(b), e(b')). For the cross-orbit comparison, for each base model b we select a different model c from the zoo in the same test-accuracy decile. Cross-orbit similarity is cos(e(b), e(c)).

We define the invariance gap as:

gap = mean(within-orbit similarity) $-$ mean(cross-orbit same-property similarity)

A positive gap indicates NFT embeds functionally identical models as less similar than property-matched models --- the encoder is non-invariant to that symmetry and wastes capacity on the symmetry-induced variation. A negative gap indicates the encoder already embeds orbit members as more similar than cross-orbit property-matched pairs --- approximate invariance. Bootstrap 95\% CIs (n\_boot=1,000) are computed on the gap.

\subsubsection{3.4 Canonicalization Conditions}

We evaluate five experimental conditions:

\begin{table}[htbp]
\centering
\resizebox{\columnwidth}{!}{%
\begin{tabular}{lll}
\toprule
Condition & Preprocessing & Purpose \\
\midrule
A & Raw weights (no preprocessing) & NFT baseline \\
B & Scaling canonicalization only & Isolate scaling effect \\
C & Sign-flip canonicalization only & Isolate sign-flip effect \\
D & Scaling + sign-flip canonicalization & Full canonicalization \\
E & Random normalization control & Test normalization artifact hypothesis \\
\bottomrule
\end{tabular}
}
\end{table}

\textbf{Scaling canonicalization (Condition B).} Per-layer L2 normalization: each hidden neuron's incoming weight vector is scaled to unit L2 norm, with the corresponding outgoing weights scaled by the inverse. Always unique and well-defined.

\textbf{Sign-flip canonicalization --- majority-sign (Conditions C, D).} For each hidden neuron i, compute the sign majority of row $W_{1}[$:,i]. If more weights are negative than positive, flip: $W_{1}[$:,i] $\leftarrow -W_{1}[$:,i], $W_{2}[i$,:] $\leftarrow -W_{2}[i$,:]. This requires a strict majority (no ties) for a unique canonical form.

\textbf{Random normalization control (Condition E).} Each hidden neuron's incoming weights are scaled by a random factor drawn from N(1, 0.1), preserving sign structure. If Condition D outperforms E, the effect is symmetry-specific; if E $\geq$ D, the effect is a normalization artifact.

\subsubsection{3.5 Canonicalization Uniqueness Audit}

For d\_in=784 (even), a tie occurs when exactly 392 weights in a neuron row are positive and 392 are negative. The binomial probability of this event is:

P(tie) = C(784, 392) $\times$ 0.$5^{784} \approx$ 2.8\%

For h=64 neurons per model, the probability of at least one tie per model is approximately $1 - (1 - 0.028)^{64} \approx 83\%$. We empirically audit all N=500 zoo models, recording fraction with unique canonical form, tied neuron count per model, and idempotency.

\subsubsection{3.6 Geometric Concentration Analysis}

We apply PCA to raw (Condition A) and canonicalized (Condition D) weight matrices, measuring:

\begin{itemize}
\item \textbf{Explained Variance Ratio (EVR)} at k $\in$ {10, 20, 50} principal components
\item \textbf{Linear regression $R^2$} from top-k PCs onto each property label, with bootstrap 95\% CIs (n\_boot=1,000, seed=42)
\end{itemize}

If canonicalization concentrates property-relevant information, we expect EVR to increase and $R^2$ to improve.

\subsubsection{3.7 Property Prediction Protocol}

For each condition (A--E), we train NFT from scratch on 400 training models (N=500 zoo, 80/10/10 split: 400 train, 50 validation, 50 test), using the corresponding preprocessed weights as input. Hyperparameters: Adam optimizer, lr=1e-3, weight\_decay=1e-4, max\_epochs=50, early stopping on validation Spearman $\rho$ (patience=10). Evaluation: Spearman $\rho$ on the 50-model held-out test split, averaged across 3 random seeds (42, 123, 456). Bootstrap 95\% CIs (n\_boot=1,000) on Spearman $\rho$ and condition differences ($\Delta\rho$). At n=50 test samples, CI width is approximately $\pm 0.3$, yielding CI width $\approx 0.6$ for any condition comparison.

